\section{Introduction}
\label{sec:intro}

\noteCLAUDE{Context: large-scale and screening.}\noteLE{Verified.}
Large-scale convex optimization is central to modern machine learning. In recent years, \emph{screening} methods have emerged as an effective technique for reducing problem size before or during optimization, thereby accelerating the solution of large-scale learning problems \citep{elghaoui2010}.

\noteCLAUDE{Classification of screening}
\noteLE{Verified.}
Screening methods have been studied from several complementary perspectives. According to whether they guarantee preservation of the optimal solution, they are broadly classified into safe screening methods \citep{elghaoui2010} and unsafe screening methods \citep{tibshirani2012strong,fan2018sure}. According to how screening is
performed, they may be static
\citep{elghaoui2010}, dynamic
\citep{bonnefoy2014dynamic}, or sequential
\citep{wang2013lasso}. Recent advances mainly focus on constructing
tighter safe regions, including ball
\citep{ndiaye2017gap,tran2025one}, dome
\citep{tran2022beyond}, and ellipsoid
\citep{dai2012ellipsoid} regions. The screening paradigm has also been generalized to related reduction mechanisms, including squeezing \citep{elvira2020safe}, relaxing \citep{guyard2022screen}, and peeling \citep{guyard2023safe}.

\noteCLAUDE{SCOPE}\noteLE{Verified.}
This paper is concerned with a different axis of classification,
namely the screening target: \emph{feature screening} and \emph{sample screening}.
Rather than studying the safety conditions guaranteeing that a
screening rule preserves the optimal solution, we study the
mathematical structure underlying screening itself.

\noteCLAUDE{Historical development.}\noteLE{Verified.}
Feature screening was first introduced for sparse learning models such
as the Lasso, where inactive features can be identified and removed
before optimization \citep{elghaoui2010safe}. Sample screening was
later developed for the soft-margin support vector machine
\citep{ogawa2014sample}. Historically, it was not introduced as an
independent mathematical operation. Instead, the feature-screening
philosophy was transferred to the SVM's dual problem, where the
dual variables correspond to training samples, and the resulting
elimination rule was interpreted as sample screening. Since then,
feature screening and sample screening have evolved largely in
parallel. They are now combined in simultaneous screening methods
\citep{shibagaki2016simultaneous}, and several recent frameworks derive
both through Fenchel--Rockafellar (FR) duality
\citep{ndiaye2017gap,yamada2021dynamic,tran2025one}.

\noteCLAUDE{Tension and gap.}\noteLE{Verified.}
This historical development suggests a close relationship between feature screening and sample screening. Yet this relationship has never been formulated mathematically. Existing work derives screening rules for particular optimization models rather than defining screening itself as a model-independent mathematical transformation. Consequently, although feature screening and sample screening are widely regarded as primal--dual counterparts, the duality between them remains an informal intuition rather than a precise mathematical statement.

\begin{center}
\emph{This naturally raises two fundamental questions. Can feature
screening and sample screening be formulated as mathematical
transformations independent of specific optimization models? If so, in
what mathematical sense are they dual?}
\end{center}


\noteCLAUDE{Central idea.}\noteLE{Verified.}
We answer these questions by formulating screening as
model-independent transformations on FR representations. Defining
screening independently of any particular optimization model requires
a common class of optimization problems on which every screening
procedure acts. Once such a class is identified, both
FR duality and screening become transformations on
the same mathematical objects, making their relationship meaningful
and directly comparable.

\noteCLAUDE{Class of problems.}\noteLE{Verified.}
Various screening methods act on convex optimization problems of the form
\[
    \min_{x\in\mathbb R^n}
    p(x)=f(Ax)+g(x)+a,
    \qquad
    A\in\mathbb R^{m\times n},
    \qquad
    a\in\mathbb R,
\]
with separable functions \(f\) and \(g\). The scalar component \(a\) is included to ensure that this class of representations is stable under the transformations considered in this paper.
In machine learning, the columns of \(A\) represent features, whereas
its rows represent samples.
This class encompasses many standard learning models, including the Lasso, Elastic Net \citep{guyard2022screen}, logistic regression \citep{wang2014safe}, Huber regression \citep{chen2020safe}, Kullback--Leibler regression \citep{dantas2021safe}, and support vector machines (SVMs) \citep{zhai2020safe,nguyen2026gap}.



\noteCLAUDE{Class of objects.}\noteLE{Verified.}
We encode every optimization problem as a quadruple
\(
p=(f,g,A,a),
\)
called a \emph{FR representation}, and denote by
\(\Gamma\) the class of all FR representations. Since \(\Gamma\) is
closed under FR duality, it provides a common
mathematical language in which feature screening and sample screening
can be formulated as transformations, rather than procedures tied to
particular optimization models.

Feature screening, restated and generalized in this paper, assumes
that a block of variables satisfies \(x_{\bar I}=z\), where
\(I\sqcup\bar I=\{1,\ldots,n\}\), and transforms the representation
into\footnote{Throughout the introduction we omit the subscript
parameters of \(S^F\) and \(S^S\) for readability.}
\[
    S^F(p)
    =
    f(A_Ix_I+A_{\bar I}z)
    +g_I(x_I)
    +a
    +g_{\bar I}(z).
\]
The classical feature screening is recovered by the special case
\(z=0\).

Sample screening, restated and generalized in this paper, assumes
\(s\in\partial f_{\bar J}(A_{\bar J}x)\), where
\(J\sqcup\bar J=\{1,\ldots,m\}\). By the Fenchel--Young equality,
\(
f_{\bar J}(A_{\bar J}x)
=
\langle s,A_{\bar J}x\rangle
-
f_{\bar J}^*(s),
\)
the representation becomes
\[
    S^S(p)
    =
    f_J(A_Jx)
    +g(x)
    +\langle A_{\bar J}^{\top}s,x\rangle
    +a
    -f_{\bar J}^*(s).
\]
The classical sample screening for SVMs is recovered by the special case
\(s=0\).

\noteCLAUDE{Structural insight.}\noteLE{Verified.}
To understand the duality between feature screening and sample screening, we seek more primitive transformations from which both can be constructed.
We show that both admit the
same structural decomposition into a translation followed by a restriction, and both primitive transformations are \emph{equivariant} under
FR duality.

\noteCLAUDE{Answer to Q2.}\noteLE{Verified.}
This immediately answers the second research question. Since every
screening operator decomposes into two equivariant primitives,
screening itself is equivariant under FR duality:
\[
    (S^F(p))^*
    =
    S^S(p^*),
\]
where \((\cdot)^*\) denotes the
FR duality. This identity is the mathematical
formulation of the duality between feature screening and sample
screening: feature screening on a primal representation is exactly
sample screening on its dual representation, expressed by the
commutative diagram
\[
\begin{array}{ccc}
p & \xrightarrow{S^F} & S^F(p)\\
{\scriptstyle *}\big\downarrow & & \big\downarrow{\scriptstyle *}\\
p^* & \xrightarrow{S^S} & S^S(p^*)
\end{array}
\]

\noteCLAUDE{Contributions.}\noteLE{Verified.}
This paper makes two main contributions. First, we formulate feature
screening and sample screening as model-independent transformations on
FR representations, providing a common mathematical language for
screening beyond individual optimization models. Second, we provide a
precise mathematical formulation of the duality between feature
screening and sample screening by proving that these transformations
are dual in the sense above. To establish these results, we
introduce the class of FR representations together with the primitive
transformations of FR duality, translation, and
restriction, and show that translation and restriction are equivariant
under FR duality, so that the screening duality
theorem follows from the structural decomposition of screening. As
direct consequences, intrinsic complexity reduction and simultaneous
screening follow naturally from this structural theory. By separating
the mathematical structure of screening from the additional
certificates required for safety, the proposed framework also provides
a foundation for future extensions to safe screening methods.


\noteCLAUDE{Organization.}\noteLE{Verified.}
The rest of the paper is organized as follows. Section~\ref{sec:operators} introduces FR representations together with the primitive transformations including duality, translation and restriction acting on them, providing the common mathematical environment for screening. Section~\ref{sec:screening} then defines feature screening and sample screening as transformations on this class of representations, proves that both decompose into translation and restriction, and establishes feature--sample duality as the main theorem. In this section, the intrinsic complexity reduction and simultaneous screening follow as immediate consequences.
